% Wersja źródłowa: audyt II, jednostka 28
\setcounter{section}{27}
\noindent\textbf{Wyniki wykorzystywane.}
Przyjmujemy twierdzenia Hahna--Banacha i Banacha--Steinhausa,
teorię baz Schaudera i norm operatorowych, podstawy przestrzeni
Hilberta oraz kryterium ośrodkowości. Wszystkie nowe konstrukcje
i ich zależności są formułowane i dowodzone w tym rozdziale.
\par\medskip
\section{Własności aproksymacji przestrzeni Banacha}
\label{sec:aproksymacja28}

W całym rozdziale $X,Y,Z$ są przestrzeniami Banacha nad
$\K\in\{\R,\C\}$. Przez $\mathcal L(X,Y)$ oznaczamy
przestrzeń operatorów liniowych ciągłych z normą operatorową, zaś
$\mathcal F(X,Y)$ -- zbiór tych $T\in\mathcal L(X,Y)$, których
obraz ma skończony wymiar. Wyrażenie \emph{operator skończonego rzędu}
nie oznacza operatora będącego jedynie odwzorowaniem algebraicznym:
zawsze wymagana jest jego ciągłość. W rozdziale o bazach Schaudera
wykazano, że rzuty częściowe takiej bazy są jednostajnie ograniczonymi
operatorami skończonego rzędu. Tutaj wyjaśnimy, jaka własność
przestrzeni za tym stoi i dlaczego nie każda przestrzeń Banacha ją ma.

\subsection{Aproksymowanie jednostajnie na zbiorach zwartych}

\begin{definicja}[AP, BAP i MAP]
\label{def:ap-bap-map28}
Niech $X$ będzie przestrzenią Banacha, a $\lambda\ge1$.
Mówimy, że $X$ ma \emph{własność aproksymacji} (AP), jeżeli dla
każdego zwartego $K\subseteq X$ i każdego $\varepsilon>0$
istnieje $T\in\mathcal F(X,X)$ spełniający
$\sup_{x\in K}\|Tx-x\|<\varepsilon$.
Mówimy, że $X$ ma \emph{$\lambda$-ograniczoną własność aproksymacji}
($\lambda$-BAP), jeżeli można ponadto zawsze wybrać taki $T$
ze $\|T\|\le\lambda$. Własność BAP oznacza istnienie
$\lambda<\infty$ spełniającego ten warunek, natomiast
\emph{metryczna własność aproksymacji} (MAP) jest równoznaczna
z $1$-BAP.
\end{definicja}

Przestrzenie skończenie wymiarowe mają MAP: wystarczy wziąć
$T=I_X$. W przestrzeniach nieskończenie wymiarowych wymaganie
zbieżności w normie operatorowej byłoby zbyt silne.

\begin{stwierdzenie}[Brak aproksymacji identyczności w normie operatorowej]
\label{prop:no-norm-identity28}
Jeżeli $X$ jest nieskończenie wymiarową przestrzenią Banacha,
to dla każdego $T\in\mathcal F(X,X)$ zachodzi
$\|I_X-T\|\ge1$.
\end{stwierdzenie}
\begin{proof}
Ponieważ $T(X)$ ma wymiar skończony, $T$ nie może być
różnowartościowy na nieskończenie wymiarowej przestrzeni $X$.
Wybieramy $x\in\ker T$ o $\|x\|=1$. Wtedy
$\|(I_X-T)x\|=1$, co daje tezę.
\end{proof}

\begin{lemat}[Kryterium BAP na zbiorach skończonych]
\label{lem:finite-bap28}
Niech $X$ będzie przestrzenią Banacha, $\lambda\ge1$.
Następujące warunki są równoważne:
\begin{enumerate}[label=\textup{(\roman*)}]
\item $X$ ma $\lambda$-BAP;
\item dla wszystkich $m\ge1$, $x_1,\ldots,x_m\in X$ i
$\varepsilon>0$ istnieje $T\in\mathcal F(X,X)$ takie,
że $\|T\|\le\lambda$ oraz
$\max_{1\le j\le m}\|Tx_j-x_j\|<\varepsilon$.
\end{enumerate}
\end{lemat}
\begin{proof}
Pierwszy warunek implikuje drugi, ponieważ zbiór skończony
jest zwarty. W przeciwną stronę ustalmy zwarty $K\subset X$
i $\varepsilon>0$. Dla $K=\varnothing$ teza jest oczywista.
W przeciwnym razie istnieją $x_1,\ldots,x_m\in K$
takie, że dla każdego $x\in K$ pewien $j$ spełnia
$\|x-x_j\|<\varepsilon/[3(1+\lambda)]$.
Wybierzmy $T$ jak w (ii) z błędem mniejszym niż
$\varepsilon/3$. Wówczas
\[
\|Tx-x\|\le\|T(x-x_j)\|+\|Tx_j-x_j\|+\|x_j-x\|
<\varepsilon.
\]
W dowodzie korzystamy jednocześnie ze skończoności siatki
$\varepsilon$ i ze \emph{wspólnej granicy norm} aproksymantów.
\end{proof}

\begin{wniosek}[Hierarchia własności]
\label{cor:hierarchy28}
Dla każdej przestrzeni Banacha zachodzi
$\mathrm{MAP}\Rightarrow\mathrm{BAP}\Rightarrow\mathrm{AP}$.
Jeśli $X$ ma $\lambda$-BAP, to ma też $\mu$-BAP
dla każdego $\mu\ge\lambda$.
\end{wniosek}
\begin{proof}
Dwie pierwsze implikacje i monotoniczność stałej wynikają
bezpośrednio z Definicji~\ref{def:ap-bap-map28}.
\end{proof}

\begin{uwaga}[Dlaczego AP nie definiuje się na zbiorach skończonych?]
\label{rem:finite-exact28}
Dla \emph{dowolnej} przestrzeni Banacha $X$ i skończonego
zbioru $\{x_1,\ldots,x_m\}\subset X$ można znaleźć
$T\in\mathcal F(X,X)$ takie, że $Tx_j=x_j$ dla każdego $j$.
Istotnie, niech $E=\operatorname{span}\{x_1,\ldots,x_m\}$
i wybierzmy bazę $u_1,\ldots,u_r$ przestrzeni $E$.
Funkcjonały współrzędnych na $E$ są ciągłe, a twierdzenie
Hahna--Banacha przedłuża je do $f_1,\ldots,f_r\in X'$.
Operator $Tx=\sum_{j=1}^rf_j(x)u_j$ jest ciągły,
skończonego rzędu i jest identycznością na $E$.
Jego norma może jednak zależeć w sposób niekontrolowany
od wyboru $E$. To właśnie odróżnia powyższy fakt od BAP.
\end{uwaga}

\subsection{Topologia zbieżności na zbiorach zwartych i granice operatorów}

\begin{definicja}[Zbieżność zwarta i silna operatorów]
\label{def:topologies-ap28}
Niech $X,Y$ będą przestrzeniami Banacha. Topologia
\emph{zbieżności jednostajnej na zbiorach zwartych}
na $\mathcal L(X,Y)$ jest wyznaczona przez seminormy
$p_K(T)=\sup_{x\in K}\|Tx\|$, gdzie $K$ przebiega
zbiory zwarte w $X$. \emph{Silna topologia operatorowa}
(SOT) jest wyznaczona przez seminormy $p_x(T)=\|Tx\|$
indeksowane przez $x\in X$. W obu przypadkach
zbieżność rozumiemy także dla sieci, a nie tylko ciągów.
\end{definicja}

\begin{twierdzenie}[Operatorowa charakterystyka AP i BAP]
\label{thm:operator-character-ap28}
Niech $X$ będzie przestrzenią Banacha i $\lambda\ge1$.
\begin{enumerate}[label=\textup{(\roman*)}]
\item $X$ ma AP wtedy i tylko wtedy, gdy
$I_X$ należy do domknięcia $\mathcal F(X,X)$
w topologii zbieżności jednostajnej na zbiorach zwartych.
\item $X$ ma $\lambda$-BAP wtedy i tylko wtedy, gdy
$I_X$ należy do domknięcia zbioru
$\{T\in\mathcal F(X,X):\|T\|\le\lambda\}$ w SOT.
Dla tego ograniczonego zbioru aproksymantów można
równoważnie użyć topologii zbieżności zwartej.
\end{enumerate}
\end{twierdzenie}
\begin{proof}
Pierwszy punkt jest dokładnym przepisaniem definicji AP
na język otoczeń w topologii $p_K$; skończone przecięcia
otoczeń odpowiadają sumie skończonej liczby zbiorów zwartych.
Pierwsza część (ii) wynika z Lematu~\ref{lem:finite-bap28},
ponieważ otoczenia w SOT wyznaczają warunki na skończenie
wielu wektorach. Na zbiorze operatorów o normach
nieprzekraczających $\lambda$ zbieżność punktowa do
$I_X$ implikuje jednostajną aproksymację na każdym
ustalonym zbiorze zwartym: stosujemy argument ze
skończoną siatką z dowodu Lematu~\ref{lem:finite-bap28}.
Implikacja odwrotna jest oczywista.
\end{proof}

\begin{stwierdzenie}[Silna granica sieci skończonego rzędu istnieje zawsze]
\label{prop:universal-sot28}
Dla \emph{każdej} przestrzeni Banacha $X$ istnieje sieć
$(T_E)_{E\in\mathfrak F(X)}$ operatorów skończonego rzędu,
która zbiega w SOT do $I_X$. Tutaj $\mathfrak F(X)$
oznacza zbiór wszystkich skończenie wymiarowych podprzestrzeni
$E\subseteq X$, skierowany inkluzją.
\end{stwierdzenie}
\begin{proof}
Dla każdego $E\in\mathfrak F(X)$ wybierzmy
ciągły rzut $T_E:X\to E$, konstruowany jak
w Uwadze~\ref{rem:finite-exact28}. Jeżeli $x\in X$,
to dla wszystkich $E\supseteq\operatorname{span}\{x\}$
zachodzi $T_Ex=x$. Oznacza to $T_Ex\to x$
w normie dla każdego $x\in X$, czyli zbieżność w SOT.
\end{proof}

\begin{uwaga}[Niebezpieczne utożsamienie AP ze zbieżnością punktową]
\label{rem:net-vs-ap28}
Twierdzenie~\ref{prop:universal-sot28} nie daje AP:
konstruowana sieć może nie mieć \emph{wspólnego ograniczenia}
na normy operatorowe, a nawet nie musi być punktowo
ograniczona jako cała rodzina. Zbieżność sieci nie implikuje
ograniczoności całego jej zbioru wartości. Dlatego
zastosowanie do takiej sieci twierdzenia Banacha--Steinhausa
bez dodatkowych przesłanek byłoby błędem.
Istnieją przestrzenie Banacha bez AP, mimo że również
w nich zachodzi Twierdzenie~\ref{prop:universal-sot28}.
\end{uwaga}

\begin{twierdzenie}[Ciągi aproksymantów i ośrodkowość]
\label{thm:sequence-bap28}
Niech $X$ będzie ośrodkową przestrzenią Banacha.
Wówczas $X$ ma BAP wtedy i tylko wtedy, gdy istnieje
\emph{ciąg} $(T_n)_{n\ge1}\subset\mathcal F(X,X)$
taki, że $T_nx\to x$ dla każdego $x\in X$.
Każda przestrzeń Banacha, dla której istnieje taki ciąg,
jest ośrodkowa, bez żadnego założenia ośrodkowości z góry.
\end{twierdzenie}
\begin{proof}
Załóżmy $\lambda$-BAP i ośrodkowość.
Niech $(x_j)_{j\ge1}$ będzie gęstym ciągiem w $X$.
Na mocy Lematu~\ref{lem:finite-bap28} dla każdego $n$
istnieje $T_n\in\mathcal F(X,X)$ o normie nie większej
niż $\lambda$ i takim, że $\|T_nx_j-x_j\|<1/n$
dla $1\le j\le n$. Zatem $T_nx_j\to x_j$ dla każdego $j$.
Dla dowolnego $x\in X$ i $x_j$ bliskiego $x$ mamy
$\|T_nx-x\|\le(\lambda+1)\|x-x_j\|+\|T_nx_j-x_j\|$,
co daje $T_nx\to x$.
Odwrotnie, jeżeli $T_nx\to x$ dla każdego $x$, to rodzina
$\{T_n:n\ge1\}$ jest punktowo ograniczona. Twierdzenie
Banacha--Steinhausa daje $\sup_n\|T_n\|<\infty$.
Przybliżenie na skończonych zbiorach, a następnie
Lemat~\ref{lem:finite-bap28}, zapewniają BAP.
Ponadto $\bigcup_nT_n(X)$ leży w przeliczalnej sumie
skończenie wymiarowych podprzestrzeni. Jej liniowa otoczka
ma przeliczalny zbiór gęsty, a zawiera gęsty zbiór
$\{\lim_nT_nx:x\in X\}=X$ w swoim domknięciu.
Wobec tego $X$ jest ośrodkowa.
\end{proof}

\subsection{Bazy Schaudera a metryczna własność aproksymacji}

\begin{twierdzenie}[Baza Schaudera implikuje BAP]
\label{thm:basis-implies-bap28}
Niech $X$ będzie przestrzenią Banacha z bazą Schaudera
$(e_n)_{n\ge1}$ oraz rzutami częściowymi $P_N$.
Jeżeli $K_e=\sup_N\|P_N\|$, to $X$ ma $K_e$-BAP.
W szczególności baza monotoniczna daje MAP.
\end{twierdzenie}
\begin{proof}
Każdy $P_N$ ma skończony rząd i spełnia
$\|P_N\|\le K_e$, a definicja bazy daje
$P_Nx\to x$ w normie dla każdego $x\in X$.
Lemat~\ref{lem:finite-bap28} wraz z oszacowaniem
ze skończoną siatką dowodzi $K_e$-BAP.
Jeżeli baza jest monotoniczna, $K_e=1$.
\end{proof}

\begin{przyklad}[Przestrzenie ciągowe]
\label{ex:sequence-map28}
Przestrzenie Banacha $c_0$ i $\ell^p$ dla
$1\le p<\infty$ mają MAP: rzuty obcinające ciąg
po $N$ wyrazach mają normę $1$ i zbieżność punktową
do identyczności udowodniono już w rozdziale o bazach Schaudera.
Nie powtarzamy tamtego dowodu. Inaczej wygląda sytuacja
w przestrzeni $\ell^\infty$, mimo że i ona ma MAP.
\end{przyklad}

\begin{twierdzenie}[MAP w dowolnej przestrzeni Hilberta]
\label{thm:hilbert-map28}
Każda rzeczywista lub zespolona przestrzeń Hilberta $H$,
bez założenia ośrodkowości, ma MAP.
\end{twierdzenie}
\begin{proof}
Niech $K\subseteq H$ będzie zwarty i $\varepsilon>0$.
Dla $K=\varnothing$ nie ma nic do dowodzenia.
W przeciwnym razie wybieramy skończoną $\varepsilon/2$-siatkę
$\{x_1,\ldots,x_m\}\subset K$ i oznaczamy
$E=\operatorname{span}\{x_1,\ldots,x_m\}$.
Niech $P_E$ będzie rzutem ortogonalnym na $E$.
Jeśli $E\ne\{0\}$, to $\|P_E\|=1$, a w przeciwnym
razie $\|P_E\|=0$. Ponieważ $P_Ex_j=x_j$,
dla $x\in K$ i odpowiedniego $j$ zachodzi
$\|x-P_Ex\|\le\|x-x_j\|+\|P_E(x_j-x)\|
<\varepsilon$. Operator $P_E$ ma skończony rząd.
\end{proof}

\begin{twierdzenie}[MAP przestrzeni funkcji ciągłych na odcinku]
\label{thm:c01-map28}
Przestrzeń Banacha $C([0,1],\K)$ z normą supremum
ma MAP. Istnieje ciąg operatorów skończonego rzędu
$L_n$ o normach równych $1$, zbieżny punktowo do
identyczności.
\end{twierdzenie}
\begin{proof}
Dla $n\ge1$ przyjmijmy $t_k=k/n$ ($0\le k\le n$).
W szczególności operator $L_n$ zależy liniowo od wartości
$f(t_0),\ldots,f(t_n)$, co jawnie zapewnia skończoność jego rzędu.
Niech $L_nf$ będzie funkcją ciągłą, liniową na każdym
$[t_k,t_{k+1}]$ i spełniającą $L_nf(t_k)=f(t_k)$.
Operator $L_n$ jest liniowy; jego obraz leży w
$(n+1)$-wymiarowej przestrzeni funkcji odcinkami liniowych.
Ponieważ $L_nf(t)$ jest kombinacją wypukłą
$f(t_k)$ i $f(t_{k+1})$, mamy
$\|L_nf\|_\infty\le\|f\|_\infty$. Dla funkcji stałej
$1$ zachodzi $L_n1=1$, więc $\|L_n\|=1$.
Z jednostajnej ciągłości $f$ na $[0,1]$ wynika
$\|L_nf-f\|_\infty\to0$. Stosujemy
Twierdzenie~\ref{thm:sequence-bap28} lub bezpośrednio
Lemat~\ref{lem:finite-bap28} z $\lambda=1$.
\end{proof}

\begin{twierdzenie}[MAP bez bazy Schaudera]
\label{thm:linfty-map28}
Przestrzeń Banacha $\ell^\infty(\N,\K)$ ma MAP,
lecz nie ma żadnej bazy Schaudera.
\end{twierdzenie}
\begin{proof}
Ustalmy $x^{(1)},\ldots,x^{(m)}\in\ell^\infty$ oraz
$\varepsilon>0$. Zbiór
$A=\{(x^{(1)}_n,\ldots,x^{(m)}_n):n\in\N\}\subset\K^m$
jest ograniczony, a więc ma skończoną
$\varepsilon/2$-siatkę. Przypisujemy każdemu indeksowi $n$ pierwszy z wybranych
punktów siatki leżący w odległości mniejszej niż $\varepsilon/2$
od odpowiadającego mu wektora. Otrzymujemy podział $\N$
na skończenie wiele niepustych zbiorów $A_1,\ldots,A_r$ tak,
że dla $n,k\in A_j$ zachodzi
$|x_n^{(s)}-x_k^{(s)}|<\varepsilon$
dla każdego $1\le s\le m$ (na mocy nierówności trójkąta).
Wybierzmy po jednym $n_j\in A_j$ i połóżmy
$(Ty)_n=y_{n_j}$, jeżeli $n\in A_j$.
Jest to liniowy operator skończonego rzędu,
$\|T\|=1$ i
$\|Tx^{(s)}-x^{(s)}\|_\infty<\varepsilon$ dla każdego $s$.
Lemat~\ref{lem:finite-bap28} dowodzi MAP.

Przestrzeń $\ell^\infty$ nie jest ośrodkowa:
nieprzeliczalna rodzina ciągów $\chi_A=(\chi_A(n))_{n\ge1}$,
gdzie $A$ przebiega wszystkie podzbiory $\N$,
ma parami odległość $1$ w normie supremum.
Tymczasem każda przestrzeń Banacha z bazą Schaudera
jest ośrodkowa. Zatem $\ell^\infty$ nie ma takiej bazy.
\end{proof}

\begin{uwaga}[Brak uniwersalnego ciągu w przestrzeni nieośrodkowej]
\label{rem:linfty-sequence28}
Twierdzenia~\ref{thm:linfty-map28} oraz
\ref{thm:sequence-bap28} pokazują jednocześnie, że
w $\ell^\infty$ nie istnieje \emph{ciąg} operatorów
skończonego rzędu zbieżny silnie do $I$.
Istnieje natomiast jednostajnie ograniczona \emph{sieć}
takich operatorów zbieżna na wszystkich zbiorach zwartych.
Jest to powód, dla którego definicje AP i BAP nie powinny
opierać się wyłącznie na ciągach.
\end{uwaga}

\subsection{Aproksymowanie operatorów zwartych}

\begin{definicja}[Operatory zwarte i aproksymowalne]
\label{def:compact-approximable28}
Niech $X,Y$ będą przestrzeniami Banacha.
Operator $T\in\mathcal L(X,Y)$ jest \emph{zwarty},
jeżeli $T(B_X)$ ma zwarte domknięcie w $Y$, gdzie
$B_X=\{x\in X:\|x\|\le1\}$.
Oznaczamy zbiór takich operatorów przez
$\mathcal K(X,Y)$. Operator jest \emph{aproksymowalny},
jeżeli należy do domknięcia $\mathcal F(X,Y)$
w normie operatorowej; ich zbiór oznaczamy
$\mathcal A(X,Y)$.
\end{definicja}

\begin{twierdzenie}[Operatory aproksymowalne są zwarte]
\label{thm:approximable-compact28}
Dla wszystkich przestrzeni Banacha $X,Y$ zachodzi
$\mathcal F(X,Y)\subseteq\mathcal A(X,Y)\subseteq
\mathcal K(X,Y)\subseteq\mathcal L(X,Y)$.
Ponadto $\mathcal K(X,Y)$ jest domkniętą podprzestrzenią
liniową $\mathcal L(X,Y)$.
\end{twierdzenie}
\begin{proof}
Każdy operator skończonego rzędu jest zwarty, ponieważ
jego ograniczony obraz kuli jednostkowej ma zwarte
domknięcie w skończenie wymiarowej przestrzeni.
Jeżeli $T_n\to T$ w normie operatorowej i każdy $T_n$
jest zwarty, to dla $\varepsilon>0$ wybieramy $n$ z
$\|T-T_n\|<\varepsilon/3$. Zbiór $T_n(B_X)$ jest
całkowicie ograniczony, więc pokrywają go skończenie
wiele kul promienia $\varepsilon/3$.
Otrzymujemy skończone pokrycie $T(B_X)$ kulami
promienia $\varepsilon$, a zatem całkowitą ograniczoność.
Ponieważ $Y$ jest zupełna, $\overline{T(B_X)}$ jest zwarty.
To pokazuje domkniętość zbioru operatorów zwartych.
Liniowość wynika ze zwartości sum i skalarnych
wielokrotności zbiorów względnie zwartych.
\end{proof}

\begin{twierdzenie}[AP w przestrzeni wartości i aproksymacja operatorów]
\label{thm:ap-compact-operators28}
Niech $X,Y$ będą przestrzeniami Banacha, przy czym
$Y$ ma AP. Wówczas każdy operator zwarty
$T\in\mathcal K(X,Y)$ jest aproksymowalny, tzn.
$\mathcal K(X,Y)=\mathcal A(X,Y)$.
Jeżeli $Y$ ma $\lambda$-BAP, to dla każdego
$T\in\mathcal K(X,Y)$ i każdego $\varepsilon>0$
istnieje $F\in\mathcal F(X,Y)$ z
$\|F-T\|<\varepsilon$ oraz $\|F\|\le\lambda\|T\|$.
\end{twierdzenie}
\begin{proof}
Niech $K=\overline{T(B_X)}$; z założenia jest zwarty.
Z AP przestrzeni $Y$ otrzymujemy dla każdego
$\varepsilon>0$ operator skończonego rzędu
$S:Y\to Y$ taki, że
$\sup_{y\in K}\|Sy-y\|<\varepsilon$.
Operator $F=S\circ T$ ma skończony rząd i
$\|F-T\|=\sup_{x\in B_X}\|S(Tx)-Tx\|<\varepsilon$.
Przy $\lambda$-BAP można ponadto wybrać $\|S\|\le\lambda$,
stąd $\|F\|\le\lambda\|T\|$.
\end{proof}

\begin{uwaga}[Właściwy kierunek w Twierdzeniu~\ref{thm:ap-compact-operators28}]
\label{rem:range-not-domain28}
Założenie AP dotyczy tutaj \emph{przestrzeni wartości $Y$},
a nie automatycznie przestrzeni wyjściowej $X$:
aproksymujemy zwarty zbiór $T(B_X)\subseteq Y$.
Utożsamienie tych dwóch ról bez dodatkowego twierdzenia
o dualności byłoby nieuzasadnione.
\end{uwaga}

\subsection{Trwałość własności i ograniczenia twierdzeń}

\begin{twierdzenie}[Trwałość dla podprzestrzeni dopełnionych]
\label{thm:complemented28}
Niech $X$ będzie przestrzenią Banacha i $Y\subseteq X$
domkniętą podprzestrzenią. Załóżmy, że istnieje
ciągły rzut $Q:X\to Y$, tj. $Q|_Y=I_Y$.
Jeżeli $X$ ma AP, to $Y$ ma AP.
Jeżeli $X$ ma $\lambda$-BAP, to $Y$ ma
$\mu$-BAP dla $\mu=\max\{1,\lambda\|Q\|\}$.
W szczególności, gdy $X$ ma MAP i $\|Q\|=1$,
to $Y$ ma MAP.
\end{twierdzenie}
\begin{proof}
Ustalmy zwarty $K\subseteq Y$. Jest on również
zwarty w $X$. Wybierzmy operator skończonego rzędu
$T:X\to X$ przybliżający identyczność jednostajnie
na $K$ i połóżmy $S=Q\circ T|_Y:Y\to Y$.
Operator $S$ ma skończony rząd, a dla $y\in K$
$\|Sy-y\|=\|Q(Ty-y)\|\le\|Q\|\,\|Ty-y\|$.
Ponieważ dokładność aproksymacji $T$ możemy wybrać
dowolnie dużą, otrzymujemy AP na $Y$.
Przy $\lambda$-BAP zachodzi dodatkowo
$\|S\|\le\|Q\|\lambda$, co dowodzi drugiej części.
\end{proof}

\begin{twierdzenie}[Niezmienniczość AP i BAP przy izomorfizmach]
\label{thm:isomorphic28}
Niech $X,Y$ będą przestrzeniami Banacha i niech
$U:X\to Y$ będzie liniowym izomorfizmem topologicznym.
Wtedy $X$ ma AP wtedy i tylko wtedy, gdy $Y$ ma AP.
Jeżeli $X$ ma $\lambda$-BAP, to $Y$ ma
$\mu$-BAP z $\mu=\max\{1,\lambda\|U\|\|U^{-1}\|\}$.
W szczególności istnienie BAP jest niezmiennicze
względem równoważnych norm, natomiast liczba $\lambda$
może się zmienić.
\end{twierdzenie}
\begin{proof}
Dla zwartego $K\subseteq Y$ jego przeciwobraz
$U^{-1}(K)$ jest zwarty w $X$. Niech
$T\in\mathcal F(X,X)$ przybliża identyczność na
$U^{-1}(K)$. Operator $S=UTU^{-1}\in\mathcal F(Y,Y)$
spełnia $\|Sy-y\|\le\|U\|\,\|TU^{-1}y-U^{-1}y\|$.
W przypadku BAP mamy dodatkowo
$\|S\|\le\|U\|\|U^{-1}\|\lambda$.
Argument jest symetryczny względem $X,Y$.
\end{proof}

\begin{uwaga}[Granice implikacji: kontrprzykłady głębokie]
\label{rem:enflo-figiel28}
Nie każda ośrodkowa, a nawet nie każda \emph{refleksywna}
przestrzeń Banacha ma AP: istnieje taki kontrprzykład
Enflo. Istnieją też przestrzenie Banacha mające AP,
ale niemające BAP (konstrukcja Figiela--Johnsona),
oraz mające BAP, lecz nie MAP. Klasyczne źródła
dwóch pierwszych kontrprzykładów to prace Enflo
(\emph{Acta Mathematica} 130, 1973, s.~309--317) oraz
Figiela i Johnsona (\emph{Proceedings of the AMS} 41,
1973, s.~197--200). Wyniki te \emph{przytaczamy bez dowodu}, ponieważ szczegółowe
konstrukcje wykraczają poza zakres rozdziału.
W przeciwieństwie do tych rezultatów kontrprzykład
``MAP nie implikuje istnienia bazy Schaudera'' z
Twierdzenia~\ref{thm:linfty-map28} jest całkowicie
elementarny i został udowodniony powyżej.
\end{uwaga}

\begin{uwaga}[Co naprawdę daje istnienie bazy?]
\label{rem:recap28}
Zestawiając Twierdzenia~\ref{thm:basis-implies-bap28}
i~\ref{thm:linfty-map28}, otrzymujemy ściśle
jednokierunkową zależność między własnościami bazowymi
i aproksymacyjnymi. Z drugiej strony nie wolno
zastępować \emph{jednostajnej} aproksymacji na
zbiorach zwartych dowolną aproksymacją punktową przez
sieć operatorów skończonego rzędu: ta ostatnia zachodzi
w każdej przestrzeni Banacha na mocy
Twierdzenia~\ref{prop:universal-sot28}.
\end{uwaga}
